%%%%%%
%
% $Autor: Srikanth Nanda$
% $Datum: 2020-01-18 11:15:45Z $
% $Pfad: githubtemplate/Template/Presentations/Template/slides/rename.tex $
% $Version: 4620 $
%
%
% !TeX encoding = utf8
% !TeX root = Rename
%
%%%%%%

%%%%%%%%%%%% FASTSLAM %%%%%%%%%%%%%%%
%%%%%%%%%%%%%%%%%%%%%%%%%%%%%%%%%%%%%

\Mysection{FastSLAM}

\STANDARD{Idea}{
Based on the AAAI-02 paper, FastSLAM introduces a \textbf{factored representation
of the SLAM (Simultaneous Localization and Mapping) posterior}:

\[
p(x_{1:t}, m \mid z_{1:t}, u_{1:t})
= p(x_{1:t} \mid z_{1:t}, u_{1:t})
\prod_{k} p(m_k \mid x_{1:t}, z_{1:t})
\]

This exploits the conditional independence between landmarks and robot
trajectory. The algorithm combines a \textbf{particle filter for robot paths}
with \textbf{independent EKFs for landmarks}. A tree-based structure enables
\textbf{logarithmic map updates}, achieving scalability to tens of thousands
of features.\cite{Montemerlo:2002}
}

\STANDARD{Conceptual Structure}{
\begin{center}
\begin{tikzpicture}
  \node (pf) [rectangle, draw, minimum width=3.5cm] {Particle Filter (Path)};
  \node (l1) [circle, draw, below left=1.2cm of pf] {EKF};
  \node (l2) [circle, draw, below=1.2cm of pf] {EKF};
  \node (l3) [circle, draw, below right=1.2cm of pf] {EKF};
  \draw[<->] (pf) -- (l1);
  \draw[<->] (pf) -- (l2);
  \draw[<->] (pf) -- (l3);
\end{tikzpicture}
\end{center}
}

\STANDARD{Hyperparameters}{
\begin{itemize}
  \item Number of particles
  \item Motion and measurement noise
  \item Landmark initialization covariance
  \item Tree resolution for map updates\cite{Montemerlo:2002}
\end{itemize}
}

\STANDARD{Applications}{
Large-scale outdoor SLAM, robot navigation in high-feature environments.\cite{Montemerlo:2002}
}

%%%%%%%%%%%%%%%%%%%%%%%%%%%%%%%%%%%%%
%%%%%%%%%%%% ACTIVE SLAM %%%%%%%%%%%%
%%%%%%%%%%%%%%%%%%%%%%%%%%%%%%%%%%%%%

\Mysection{Active SLAM}

\STANDARD{Idea}{
From the IEEE/ASME T-Mech paper, Active SLAM frames SLAM as a
\textbf{dual-mode MPC problem}:

\textbf{Mode 1: SLAM-uncertainty reduction}\\
Uses convex relaxation, graph topology, and Fisher Information eigenvalue
bounds to approximate future uncertainty. \cite{Chen:2018}

\textbf{Mode 2: Coverage}\\
Sequential quadratic programming ensures environment coverage while
maintaining safe trajectories.

The authors employ \textbf{Linear SLAM} (submap joining) for fast covariance
recovery and scalable planning. \cite{Chen:2018}
}

\STANDARD{Conceptual Structure}{
\begin{center}
\begin{tikzpicture}
  \node (robot) [circle, draw] {Robot};
  \node (mpc) [rectangle, draw, right=2.5cm of robot] {MPC Modes};
  \node (info) [rectangle, draw, above=1.2cm of mpc] {Uncertainty Model};
  \node (cov)  [rectangle, draw, below=1.2cm of mpc] {Coverage Model};
  \node (sub)  [rectangle, draw, right=2.5cm of mpc] {Submap Joining};
  \draw[->] (robot) -- (mpc);
  \draw[->] (mpc) -- (info);
  \draw[->] (mpc) -- (cov);
  \draw[->] (mpc) -- (sub);
\end{tikzpicture}
\end{center}
}

\STANDARD{Hyperparameters}{
\begin{itemize}
  \item MPC horizon
  \item Convex relaxation parameters
  \item SLAM measurement noise
  \item Control bounds \cite{Chen:2018}
\end{itemize}
}

\STANDARD{Applications}{
Autonomous UAV exploration, obstacle-aware navigation, area scanning. \cite{Chen:2018}
}

%%%%%%%%%%%%%%%%%%%%%%%%%%%%%%%%%%%%%
%%%%%%%%%%%% MONOSLAM %%%%%%%%%%%%%%%
%%%%%%%%%%%%%%%%%%%%%%%%%%%%%%%%%%%%%

\Mysection{MonoSLAM }

\STANDARD{Idea}{
The PAMI 2007 paper presents the first real-time monocular SLAM with drift-free
performance. MonoSLAM uses:

\begin{itemize}
  \item \textbf{EKF joint state}: camera pose + 3D features
  \item \textbf{Active measurement search}: predicts feature uncertainty ellipses
  \item \textbf{General 3D motion model} for smooth tracking
  \item \textbf{Inverse-depth feature initialization}
\end{itemize}

The emphasis is on robust localization for tasks such as humanoid robot
control and augmented reality.\cite{Davison:2007}
}

\STANDARD{Conceptual Structure}{
\begin{center}
\begin{tikzpicture}
  \node (cam)  [circle, draw] {Camera};
  \node (track) [rectangle, draw, right=2.5cm of cam] {Feature Tracking};
  \node (ekf)  [rectangle, draw, right=2.5cm of track] {EKF State};
  \node (map)  [rectangle, draw, below=1.3cm of ekf] {Persistent Map};
  \draw[->] (cam) -- (track);
  \draw[->] (track) -- (ekf);
  \draw[->] (track) -- (map);
\end{tikzpicture}
\end{center}
}

\STANDARD{Hyperparameters}{
Process noise, measurement noise, feature search windows, initialization depth.\cite{Davison:2007}
}

\STANDARD{Applications}{
Augmented reality, humanoid robots, embedded monocular SLAM.\cite{Davison:2007}
}

%%%%%%%%%%%%%%%%%%%%%%%%%%%%%%%%%%%%%
%%%%%%%%%%%% ORB-SLAM %%%%%%%%%%%%%%%
%%%%%%%%%%%%%%%%%%%%%%%%%%%%%%%%%%%%%

\Mysection{ORB-SLAM }

\STANDARD{Idea}{
According to the IEEE Transactions on Robotics paper, ORB-SLAM is designed
for real-time, robust large-scale operation:

\begin{itemize}
  \item \textbf{ORB features} for all SLAM tasks
  \item \textbf{Covisibility graph} for local mapping
  \item \textbf{Essential graph} for global optimization
  \item \textbf{Automatic initialization} via model selection
  \item \textbf{Survival-of-the-fittest} keyframe policy \cite{MurArtal:2015}
\end{itemize}

It achieves state-of-the-art accuracy across 27 sequences.
}

\STANDARD{Conceptual Structure}{
\begin{center}
\begin{tikzpicture}
  \node (track) [rectangle, draw] {Tracking};
  \node (local) [rectangle, draw, right=2.0cm of track] {Local Mapping};
  \node (loop)  [rectangle, draw, right=2.0cm of local] {Loop Closing};
  \node (graph) [rectangle, draw, below=1.3cm of local] {Essential Graph};
  \draw[->] (track) -- (local);
  \draw[->] (local) -- (loop);
  \draw[->] (local) -- (graph);
\end{tikzpicture}
\end{center}
}

\STANDARD{Hyperparameters}{
Number of ORB features, BA window size, keyframe thresholds, loop closure limits. \cite{MurArtal:2015}
}

\STANDARD{Applications}{
AR/VR localization, robotics, drone navigation, lifelong SLAM. \cite{MurArtal:2015}
}

%%%%%%%%%%%%%%%%%%%%%%%%%%%%%%%%%%%%%
%%%%%%%%%%% VISUAL-SLAM SURVEYS %%%%
%%%%%%%%%%%%%%%%%%%%%%%%%%%%%%%%%%%%%

\Mysection{Visual SLAM Surveys }

\STANDARD{Idea}{
Document insights reveal three dominant categories:

\begin{itemize}
  \item \textbf{Feature-based SLAM} (ORB-SLAM, MonoSLAM)
  \item \textbf{Direct SLAM} (LSD-SLAM, DSO)
  \item \textbf{RGB-D SLAM} (ElasticFusion, Kintinuous) \cite{MacarioBarros:2022, Taketomi:2017}
\end{itemize}

Major evaluation dimensions include:
tracking robustness, drift handling, loop closure, computational efficiency. \cite{Taketomi:2017}
}

\STANDARD{Conceptual Structure}{
\begin{center}
\begin{tikzpicture}
  \node (f) [circle, draw] {Feature};
  \node (d) [circle, draw, right=2.3cm of f] {Direct};
  \node (r) [circle, draw, right=2.3cm of d] {RGB-D};
  \draw[->] (f) -- (d) -- (r);
  
\end{tikzpicture}
\end{center}
}

%%%%%%%%%%%%%%%%%%%%%%%%%%%%%%%%%%%%%
%%%%%%%%%%%% SLAM OVERVIEW %%%%%%%%%
%%%%%%%%%%%%%%%%%%%%%%%%%%%%%%%%%%%%%

\Mysection{SLAM Overview }

\STANDARD{Idea}{
The overview document explains SLAM in terms of:

\begin{itemize}
  \item \textbf{Front-end}: feature extraction, scan matching, data association
  \item \textbf{Back-end}: optimization (EKF, smoothing, pose graphs)
  \item \textbf{Sensors}: LiDAR, camera, IMU, GNSS integration
  \item \textbf{Failure modes}: drift, wrong associations, poor initialization \cite{Alsadik:2021}
\end{itemize}
}

\STANDARD{Conceptual Structure}{
\begin{center}
\begin{tikzpicture}
  \node (front) [rectangle, draw] {Front-End};
  \node (back)  [rectangle, draw, right=3.2cm of front] {Back-End Optimization};
  \draw[<->] (front) -- (back);
\end{tikzpicture}
\end{center}
}
%%%%%%%%%%%%%%%%%%%%%%%%%%%%%%%%%%%%%%%%%%%%%%%%%%%%%%%%%%%%
\Mysection{SLAM Algorithms Code Implementation}

\STANDARD{Python Implementation SLAM Algorithms}{
\lstinputlisting[
  language=Python,
  style=mystyle,
  firstline=1,
  caption={Python Implementation SLAM Algorithms}
]{Code/SLAM.py}
}



